% ৪.৬ ওয়েবসাইটের কাঠামো (বিস্তৃত): Linear, Tree, Web-linked, Hybrid
\section{ওয়েবসাইটের কাঠামো (বিস্তৃত)}

\begin{learningbox}
এই টপিক শেষে শিক্ষার্থী website structure-এর প্রধান ধরনগুলো (Linear, Tree, Web-linked/Network, Hybrid) চেনবে, প্রতিটির ব্যবহার, সুবিধা ও অসুবিধা এবং ব্যবহারিক উদাহরণ ব্যাখ্যা করতে পারবে।
\end{learningbox}

\subsection{সংজ্ঞা ও প্রয়োজনীয়তা}
ওয়েবসাইটের কাঠামো হলো page গুলো কিভাবে সাজানো হবে এবং ব্যবহারকারী কিভাবে navigation করে target content-এ পৌঁছাবে—এর পরিকল্পিত বিন্যাস। একটি পরিষ্কার structure user experience ও searchability উন্নত করে।

\subsection{Linear Structure}
\begin{definitionbox}{Linear Structure}
Linear Structure-এ page গুলো ধারাবাহিকভাবে সাজানো থাকে—user নির্দিষ্ট ক্রমে content দেখে।\end{definitionbox}

\begin{keypointbox}{ব্যবহার ও বৈশিষ্ট্য}
\begin{itemize}
\item Step-by-step tutorial, quiz, wizard/registration flow
\item সহজ ও controlled navigation
\item সীমাবদ্ধ user freedom
\end{itemize}
\end{keypointbox}

\subsection{Tree (Hierarchical) Structure}
\begin{definitionbox}{Tree Structure}
Tree Structure-এ homepage root হিসেবে থাকে এবং main section ও subsection আকারে স্তরভিত্তিকভাবে page সাজানো হয়।\end{definitionbox}

\begin{keypointbox}{ব্যবহার ও বৈশিষ্ট্য}
\begin{itemize}
\item School/college, corporate, government websites
\item Organized categories, predictable navigation
\item ভালো for SEO এবং large content sets
\end{itemize}
\end{keypointbox}

\subsection{Web-linked / Network Structure}
\begin{definitionbox}{Web-linked Structure}
Web-linked Structure-এ page গুলো পরস্পরের সাথে বহুপথে link থাকে—user স্বাধীনভাবে related pages navigate করতে পারে।\end{definitionbox}

\begin{keypointbox}{ব্যবহার ও বৈশিষ্ট্য}
\begin{itemize}
\item Wiki, documentation, news portals
\item Rich internal linking; flexible navigation
\item বেশি link management ও broken-link risk
\end{itemize}
\end{keypointbox}

\subsection{Hybrid Structure}
\begin{definitionbox}{Hybrid Structure}
Hybrid Structure বিভিন্ন model একসাথে ব্যবহার করে—Tree core, কিন্তু tutorial sections linear ও articles web-linked হতে পারে।\end{definitionbox}

\begin{keypointbox}{ব্যবহার ও বৈশিষ্ট্য}
\begin{itemize}
\item University portals, large e-commerce, government sites
\item Flexible, scalable, but design ও maintenance জটিল
\end{itemize}
\end{keypointbox}

\subsection{Structure design tips}
\begin{itemize}
\item Home থেকে key pages সহজে পৌঁছানো যায় কিনা পরীক্ষা করো।
\item Navigation menu logical ও concise রাখো।
\item Deep nesting এড়াও; প্রয়োজন হলে breadcrumbs ব্যবহার করো।
\item Mobile users-এর জন্য responsive navigation ডিজাইন করো।
\item Sitemap তৈরি করে search engines ও users-কে সহায়তা করো।
\end{itemize}

\begin{examplebox}{College website (recommended)}
Main: Tree structure (Home → Academic, Admission, Notice, Result, Contact)\\
Admission flow: Linear (Notice → Apply → Payment → Confirmation)\\
Related articles/announcements: Web-linked for easy reference
\end{examplebox}

\begin{practicebox}
\begin{enumerate}
\item একটি ছোট website-এর জন্য Tree structure-এর sitemap আঁকো (৫ page)।
\item Tutorial section-এর জন্য Linear structure কিভাবে ব্যবহার করবে তা ব্যাখ্যা করো।
\item Web-linked structure-এ broken link সমস্যা কিভাবে হ্রাস করা যায়? তিনটি উপায় লেখো।
\end{enumerate}
\end{practicebox}
